% Generated test artifact for progress-report
% format: latex
% template: thesis-status

\documentclass[11pt,a4paper]{report}
\usepackage[margin=2.6cm]{geometry}
\usepackage{hyperref}

\title{Thesis Progress Note}
\author{Your Name}
\date{}

\begin{document}
\maketitle

\chapter*{Current Focus}
This week the main focus remained robust data augmentation for small-sample image classification, with the immediate goal of deciding whether the current direction is stable enough to justify a larger run.

\chapter*{Completed Since Last Report}

\section*{Experiments}
I re-ran the baseline with the cleaned split so the comparison now matches the evaluation setup used in the last discussion. I also aligned the current result figure, the notebook summary, and the code changes so the argument for the latest direction is easier to follow.

\section*{Writing and Notes}
I separated the discarded attempts from the retained path in the experiment notes. That cleanup makes the current story much easier to explain in a thesis-style progress note instead of a raw lab log.

\chapter*{Risk Register}
The latest result still comes from a single seed, so the gain should be treated as preliminary. GPU queue time also remains the main risk for any larger follow-up run.

\chapter*{Decisions Needed}
The main decision is whether the next step should be multi-seed verification or direct expansion to a larger dataset.

\chapter*{Next Milestones}
The next milestone is to make the current result more defensible with additional verification and then prepare a compact comparison table for the next meeting.

\end{document}
